\section{Limitations}
\label{sec:limits}

We state these in order of how much they should change a reader's confidence.

\paragraph{No human validated any screening decision.} The protocol specified an independent
human audit of a stratified sample; the human co-author declined it. The corpus is therefore
entirely AI-determined. The substitute we ran --- a third independent AI pass agreeing
\consAgree{} of the time on \nCons{} papers --- measures stability under re-application, not
correctness. Three passes drawn from one model family can share a systematic error, and the
Layer-1 disagreement structure shows this family does have systematic boundary preferences.
This is the weakness most likely to matter, and the audit worksheet remains in the repository
for anyone who wishes to perform it.

\paragraph{The coders and the screeners share a base model family.} Pass A ran on Claude
Fable 5 and Pass B on Claude Opus 5, which reduces but does not eliminate correlated error;
coding used a single pass. A review of AI research conducted by AI is exposed to exactly the
correlated-error problem \citet{bowkis2026automated} describe.

\paragraph{Abstract-based coding.} Nine of ten dimensions were coded from titles and
abstracts, with \nNeedsFulltext{} papers flagged as warranting full text. Abstracts
systematically understate implementation detail; we demonstrated the magnitude for one
dimension by measuring artifact release from full text instead, where the correction was
large (Section~\ref{sec:artifacts}). Dimensions we did not verify this way --- particularly
auditability mechanisms --- are likely to be understated in the same direction. Our
auditability figures should therefore be read as lower bounds on what these systems provide,
and the interesting quantity is not the absolute level but which papers report it at all.

\paragraph{arXiv-only, computer-science-centric.} A cross-check against OpenAlex confirmed
that flagship biological autonomous-science systems publish on bioRxiv. Our map describes the
arXiv-visible field; wet-lab-heavy work is under-represented, which probably biases our
external-validation rate downward.

\paragraph{Corpus boundaries are contested, and we chose one resolution.} The 21.5\%
disagreement rate is itself evidence that reasonable reviewers would draw this corpus
differently. We have released the boundary rules, the excluded classes with counts, and every
decision with its justification, so that a reader who disagrees can recompute the map under
different rules rather than take ours on faith.

\paragraph{Right-censoring and version drift.} The final quarter is partial, and we coded the
latest version of each paper available at harvest; papers revised after 2026-07-25 may no
longer match their codes.

\paragraph{Descriptive, not causal.} Every statistic here is a frequency. We make no claim
about whether auditability practices are improving, whether they affect research quality, or
whether any particular system's results are correct.

\section{Conclusion}

The literature on autonomous AI research systems grew roughly \growthFactor{}-fold in nine
quarters and now evaluates itself by the standards of machine-learning benchmarking:
diligently, at scale, and in a way that mostly does not let a reader reconstruct how any
particular claim was reached. \pctHeldOut{} of papers test a result outside the loop that
produced it, \pctAuditNone{} provide no auditability mechanism at all, and \pctHumanUnspec{}
do not say what humans did.

The pattern is not uniform, and its shape is the finding. Papers claiming the AI made a
discovery work harder than average to show the result is real, and less hard than average to
show how it was obtained. Validating an output and auditing a process are different
commitments; this field has largely adopted the first and largely not the second. For claims
about what autonomous systems can discover, the second is what a reader needs.

We are not in a position to be superior about this. Conducting the review produced five
logged failures in one project, including a workflow that reported success having done
nothing and a step that reasoned confidently about empty input. Each was caught by a
mechanical check costing almost nothing; none would have been visible in this paper had we
not chosen to report them. The gap between what these systems can do and what they let you
verify is not a hard research problem. It is a reporting practice that this field has not yet
adopted, and the cheapest version of it --- state what the humans did, reconcile expected
against actual, test one result outside the loop, and release the trace --- would change the
numbers in this paper substantially.

\section{Contribution statement}
\label{sec:contrib}

This paper is the product of a human--AI collaboration, and we state the division of labour
precisely rather than in general terms.

\textbf{Claude Fable 5 (Anthropic)} proposed the research direction from a survey of the
July 2026 literature; wrote the review protocol, codebook, screening prompts, and both
amendments; wrote all software (harvest, screening, coding, analysis, figures, artifact
verification, citation verification); executed the search, screening, coding, and analysis;
generated the figures; drafted the entire manuscript; and maintained the process logs,
including the record of its own errors.

\textbf{Khristian Kopachelli} initiated and framed the project; selected the research
direction from the options presented; set the publication and authorship strategy; approved
the protocol at the pre-screening checkpoint; made the scope decisions at the corpus
checkpoint, including overruling the AI's boundary rule on deep-research agents (a reversal
that changed the corpus by 183 papers); declined the specified human screening audit; and
reviewed and approved the manuscript. He takes responsibility for the work as submitted.

arXiv's policy does not permit an AI system to be listed as an author, and we comply with it.
We record here that if listing reflected contribution rather than accountability, the AI would
be a co-author of this paper. The companion deposit (Section~\ref{sec:avail}) carries a joint
authorship statement to that effect.

We note the obvious circularity --- an AI-conducted study of AI-conducted research, reporting
on its own reliability --- and refer readers to Section~\ref{sec:reflexive} for what we did to
make the claims checkable and to Section~\ref{sec:limits} for where we could not.

\section{Data, code, and artifact availability}
\label{sec:avail}

Everything needed to reproduce this paper is released: the harvest script and manifest, the
full candidate corpus, both layers of screening decisions with per-paper justifications, the
coded corpus, the artifact-verification results, the analysis and figure scripts, the protocol
with all amendments, and the three process logs. Every number in the paper is generated into a
macro file by the analysis script; none is transcribed by hand. Every reference was verified
against live arXiv and Crossref metadata by a script included in the repository.

\begin{itemize}\itemsep1pt
\item Repository: \url{https://github.com/kopachelli/arxiv-frontier}
\item Companion deposit with joint-authorship statement and archived data: (DOI on deposit)
\end{itemize}
